\documentclass[11pt,a4paper]{article}
\usepackage[T1]{fontenc}
\usepackage[utf8]{inputenc}
\usepackage{amsmath,amssymb,xcolor}
\usepackage[margin=25mm]{geometry}
\usepackage[colorlinks=true,urlcolor=blue]{hyperref}
\definecolor{nexorablue}{HTML}{2459B5}
\setlength{\parindent}{0pt}
\setlength{\parskip}{7pt}
\setlength{\emergencystretch}{2em}
\renewcommand{\familydefault}{\sfdefault}
\begin{document}

{\small\bfseries\color{nexorablue} NEXORA CPGE}\par
{\small MPSI / NEX-LOG-06 / Corrigé}\par
{\LARGE\bfseries Constituer un montant avec deux types de jetons}\par
\vspace{0.25em}\hrule\vspace{0.6em}

\section*{Quatre bases}
On réalise $18=7+7+4$, $19=7+4+4+4$, $20=5\times4$ et $21=3\times7$. Ces quatre cas couvrent les quatre classes modulo $4$.

\section*{Hérédité forte}
Pour $n\ge22$, l’entier $n-4$ appartient à $[18,n-1]$. Par l’hypothèse de récurrence forte, il est réalisable. L’ajout d’un jeton de valeur $4$ donne $n$.

\section*{Montant impossible}
Pour former $17$, on pourrait prendre zéro, un ou deux jetons de valeur $7$. Les restes seraient $17$, $10$ ou $3$, aucun ne pouvant être obtenu uniquement avec des jetons de valeur $4$.


\end{document}
